\section{Discussion}
The enhanced performance of the synapse-engineered QDSSC stems from two major factors: improved optical absorption and suppressed charge recombination. The high specific surface area of the mesoporous network, templated by the Apex-9 surfactant, allowed for a denser packing of CdSe quantum dots, which directly boosted light-harvesting capacity and increased $J_{\text{sc}}$.

Furthermore, the significant improvements in $V_{\text{oc}}$ and fill factor point to reduced carrier recombination. To confirm this, we conducted electrochemical impedance spectroscopy (EIS). The charge recombination resistance ($R_{\text{ct}}$) at the $\text{TiO}_2$/electrolyte interface was determined from the Nyquist plot. The synapse-engineered device exhibited a recombination resistance of $38.4\,\Omega$, which is more than double that of the baseline cell ($15.2\,\Omega$). This indicates that the synapse-engineered interface effectively passivates surface defect states in $\text{TiO}_2$, which act as traps and recombination centers under open-circuit conditions. These results align with previous observations on structured oxide interfaces by Thorne \& Vance \cite{thorne2025} and surfactant-assisted hydrothermal synthesis routes developed by Vance \& Sterling \cite{vance2023}. 

While the synapse-engineered cell outperforms the baseline and the recent benchmark reported by Sterling \cite{sterling2024}, stability remains an open challenge. Under continuous AM 1.5G illumination for $100\,\text{hours}$, the PCE of the synapse-engineered device declined by $12\%$. This is primarily due to photothermal degradation at the quantum dot interface. Future work will investigate the deployment of an ultrathin atomic layer deposited (ALD) ZnS passivation layer to shield the QDs and enhance stability, as suggested by the Apex Solar Consortium \cite{apexsolar}.
